\section{Technical graphical and propagation details}
\label{app:technical-background}

\subsection{Local finiteness and backward predecessors}
\label{app:local-finiteness}

Let
\begin{equation*}
\overline\alpha
=
\sup_{i\in\Lambda}
\bb{
 \sum_{S\subseteq N(i)}\delta_i(S)
 +
 \sum_{\substack{S\subseteq N(i)\\S\neq\vn}}\beta_i(S)
},
\qquad
m_0=\sup_{i\in\Lambda}\abs{N(i)}.
\end{equation*}
Both constants are finite under the standing bounded-rate and bounded-degree
assumptions.

\begin{lemma}[Local finiteness]
\label{lem:dual-local-finiteness}
If the signed dual starts from~$A\Subset\Lambda$, then almost surely only
finitely many relevant interactions occur on every bounded time interval.
Consequently,~$\mathcal I_T$ and~$\mathcal P_T$ are finite for every~$T<\infty$,
ordinary interaction times are distinct, and every backward predecessor chain of
patches whose boundary times decrease strictly reaches time~$0$ after finitely
many steps.
\end{lemma}

\begin{proof}
When the active set is~$B$, the total rate of relevant clocks is at most
$\overline\alpha\abs B$. One interaction adds at most~$m_0$ sites. Hence the
number of active sites and the number of relevant rings are dominated by a
continuous-time branching process in which each particle rings at rate
$\overline\alpha$ and creates at most~$m_0$ new particles. This branching
process is nonexplosive. Independent Poisson processes have no simultaneous
points almost surely, and a countable union over pairs of clocks preserves this
property. The remaining conclusions follow from the definitions of the
successful skeleton and the patch family.
\end{proof}

In particular,
\begin{equation}
\label{eq:finite-patch-family-last-interaction}
\cb{\abs{\mathcal P}<\infty}
=
\bigcup_{n\in\N}L_n.
\end{equation}
Indeed, a finite patch family contains only finitely many successful
interactions and therefore has a last one. Conversely, on~$L_T$ there are only
finitely many interactions before~$T$ and none after it.

\subsection{Confined-interaction identity}
\label{app:confined-interactions}

For~$R\subseteq\Lambda$, let~$\eta^{R,0}$ agree with~$\eta$ on~$R$ and vanish on
$R^c$. Define the modified flip rates
\begin{equation*}
c_{i,R}(\eta)
=
\pw{c_i\bb{\eta^{R,0}}}{i\in R}
{c_i\bb{\eta^{\cb{i},0}}}{i\notin R},
\end{equation*}
and let~$P_t^R$ be the corresponding semigroup. The process inside~$R$ is the
original spin system with zero boundary, while the sites outside~$R$ evolve
independently. Hence, if~$f$ depends only on the spins in~$R$,
\begin{equation}
\label{eq:modified-zero-boundary-agreement}
P_t^Rf=P_t^{R,0}f.
\end{equation}

\begin{proposition}[Confined-interaction identity]
\label{prop:confined-interaction-identity}
Let~$A\subseteq R\Subset\Lambda$,~$0\leq T\leq t$, and let~$\mu$ be any
probability measure. Then
\begin{equation}
\label{eq:confined-interaction-identity}
\E_A\Cb{W_t^\mu\ind\bb{E_T^R}}
=
\mu\bb{P_{t-T}P_T^{R,0}\chi_A}.
\end{equation}
\end{proposition}

\begin{proof}
For a fixed configuration~$\eta$, write
\begin{equation*}
Z_t^\eta
=
\sigma_t
\exp\bb{\int_0^tV(A_s)\,ds}
\chi_{A_t}(\eta).
\end{equation*}
The pathwise patch factorization and Theorem~\ref{thm:patch-factorization} give
$W_t^{\delta_\eta}=\CE{Z_t^\eta}{\mathcal G_t}$. Since~$E_T^R\in\mathcal G_t$,
\begin{equation}
\label{eq:confined-weight-to-dual}
\E_A\Cb{W_t^{\delta_\eta}\ind\bb{E_T^R}}
=
\E_A\Cb{Z_t^\eta\ind\bb{E_T^R}}.
\end{equation}

During the dual interval~$[0,T]$, retain every empty-target death interaction
and retain a nonempty-target split or birth only when its source lies in~$R$ and
its target is contained in~$R$. Let~$\Dcal_R$ be the resulting signed-dual
generator. If the active set is~$B$, the total rate of discarded interactions
is
\begin{equation*}
\kappa_R(B)
=
\sum_{i\in B}
\sum_{\substack{\vn\neq S\subseteq N(i)\\i\notin R\text{ or }S\nsubseteq R}}
\bb{\delta_i(S)+\beta_i(S)}.
\end{equation*}
The indicator of~$E_T^R$ kills the dual at the first discarded successful
interaction, so the Feynman--Kac operator on~$[0,T]$ is
\begin{equation*}
\Dcal_R+V-\kappa_R.
\end{equation*}

This is exactly the monomial-dual operator of the modified rates~$c_{i,R}$. For
$i\in R$, the multilinear expansion of~$c_{i,R}^x$ retains precisely the
coefficients~$c_i^x(S)$ with~$S\subseteq R$. For~$i\notin R$, only the
empty-neighbour coefficient remains. Thus the nonempty-target dual transitions
are exactly those retained in~$\Dcal_R$, every empty-target death remains, and
the discarded total rate subtracts~$\kappa_R$ from the potential.

Apply the modified duality on~$[0,T]$ and the original duality on~$[T,t]$.
Dual intervals act from right to left on observables, so
\begin{equation*}
\E_A\Cb{Z_t^\eta\ind\bb{E_T^R}}
=
\bb{P_{t-T}P_T^R\chi_A}(\eta).
\end{equation*}
Since~$A\subseteq R$,~$\chi_A$ depends only on~$R$, and
\eqref{eq:modified-zero-boundary-agreement} replaces~$P_T^R$ by
$P_T^{R,0}$. Integrating the resulting pointwise identity over~$\mu$ proves
\eqref{eq:confined-interaction-identity}.
\end{proof}

\subsection{Continuation without successful interactions}
\label{app:no-late-continuation}

\begin{proposition}[No-late continuation]
\label{prop:no-late-continuation-identity}
Fix~$0\leq T<t$. For every probability measure~$\mu$,
\begin{align}
\CE{W_t^\mu\ind\bb{L_{T,t}}}{\mathcal G_T}
&=
\prod_{P\in\mathcal B_T}C(P)\,
\mu\bb{
 \prod_{P\in\mathcal E_T}
 C\bb{\eta\bb{i(P)},P^{\uparrow t}}
},
\label{eq:appendix-finite-continuation-identity}
\\
\CE{W\ind\bb{L_T}}{\mathcal G_T}
&=
\prod_{P\in\mathcal B_T}C(P)
\prod_{P\in\mathcal E_T}C\bb{P^{\uparrow\infty}}.
\label{eq:appendix-infinite-continuation-identity}
\end{align}
The same finite-horizon identity holds with every contribution replaced by its
$\Lcal^\varepsilon$ counterpart.
\end{proposition}

\begin{proof}
For~$\eta\in\Omega$, use the random variable~$Z_t^\eta$ from the proof of
Proposition~\ref{prop:confined-interaction-identity}. The patch representation
gives~$W_t^{\delta_\eta}=\CE{Z_t^\eta}{\mathcal G_t}$, and
$L_{T,t}\in\mathcal G_t$. The tower property therefore gives
\begin{equation*}
\CE{W_t^{\delta_\eta}\ind\bb{L_{T,t}}}{\mathcal G_T}
=
\CE{Z_t^\eta\ind\bb{L_{T,t}}}{\mathcal G_T}.
\end{equation*}

Conditional on~$\mathcal G_T$, the hidden data of the bulk patches have their
consistent patch laws. The future Poisson increments based at the distinct
sites of~$\mathcal E_T$ are independent of these data and of one another. As
long as no successful interaction occurs, no site outside~$\mb{Cone}_T$ can
become dual-active. Hence~$L_{T,t}$ is the intersection, over
$P\in\mathcal E_T$, of the local events that the continuation of~$P$ creates no
successful interaction before~$t$. Patch factorization gives the bulk factors
$C(P)$. For an end patch~$P$, the one-site expectation of its Feynman--Kac
factor, terminal value~$\eta(i(P))$, and local no-success event is precisely
$C(\eta(i(P)),P^{\uparrow t})$; this is the meaning of the continuation identity
\eqref{eq:end-patch-continuation}. Multiplying these independent local
expectations proves the pointwise version of
\eqref{eq:appendix-finite-continuation-identity}. Integration over~$\eta$ with
respect to~$\mu$ proves the stated formula. The proof for~$\Lcal^\varepsilon$
is identical.

On~$L_T$, the full patch family consists of~$\mathcal B_T$ and one infinite
continuation of each patch in~$\mathcal E_T$. The latter family is finite by
Lemma~\ref{lem:dual-local-finiteness}. The same conditional product
disintegration over the independent future increments gives one local factor
$C(P^{\uparrow\infty})$ per end patch. The existence and value of each factor
are the~$u\to\infty$ limits in
\eqref{eq:end-patch-continuation}, computed explicitly in
Appendix~\ref{app:positivity-algebra}. This proves
\eqref{eq:appendix-infinite-continuation-identity}.
\end{proof}

\subsection{Finite-volume and finite-propagation approximations}
\label{app:finite-propagation}

Let~$d$ be the graph metric used in Section~\ref{sec:convergence}, and choose
$1\leq\ell<\infty$ such that~$N(i)\subseteq B(i,\ell)$ for every~$i$. Put
\begin{equation*}
\overline c
=
\sup_{i\in\Lambda,\,\eta\in\Omega}c_i(\eta),
\qquad
m
=
\sup_{i\in\Lambda}\abs{B(i,\ell)}.
\end{equation*}
Both constants are finite.

\begin{lemma}[Finite propagation]
\label{lem:finite-propagation}
Let~$f$ be supported on~$A\Subset\Lambda$. For every~$a>0$, there are
$v<\infty$ and~$C_A<\infty$ such that
\begin{equation}
\label{eq:finite-propagation}
\norm{\bb{P_t-P_t^{R,0}}f}_\infty
\leq
C_Ae^{-at}\norm f_\infty
\end{equation}
for every~$t\geq0$ and every finite~$R$ containing~$B(A,vt)$.
\end{lemma}

\begin{proof}
Construct the process from rate-$\overline c$ candidate clocks and independent
acceptance marks. Explore the clocks backwards from~$A$ at time~$t$. Whenever
a candidate clock at~$i$ is encountered, add~$B(i,\ell)$ to the set of sites
whose earlier histories must be known. If this backward influence cluster stays
inside~$R$, then the infinite-volume process and the zero-boundary process,
coupled with the same clocks and marks in~$R$, agree on~$A$ at time~$t$.

To leave~$B(A,r)$, the cluster must contain a chronological dependency path of
at least~$n=\lfloor r/\ell\rfloor+1$ candidate clocks. There are at most
$\abs A m^k$ spatial paths of length~$k$, and the expected number of
chronologically ordered~$k$-tuples of clocks along a fixed path is
$(\overline c t)^k/k!$. Therefore, for every~$\theta>0$,
\begin{align*}
\pP\bb{\text{the influence cluster leaves }B(A,r)}
&\leq
\abs A\sum_{k\geq n}\frac{\bb{m\overline c t}^k}{k!}
\\
&\leq
\abs A\exp\bb{m\overline c t e^\theta-\theta n}.
\end{align*}
Take~$r=vt$ and choose~$v$ so that
\begin{equation*}
\frac{\theta v}{\ell}-m\overline c e^\theta>a.
\end{equation*}
After increasing the constant to cover bounded~$t$, the exit probability is at
most~$C_Ae^{-at}$. The graphical coupling gives
\begin{equation*}
\abs{P_tf(\eta)-P_t^{R,0}f(\eta)}
\leq
2\norm f_\infty
\pP\bb{\text{the influence cluster leaves }R},
\end{equation*}
uniformly in~$\eta$, and proves~\eqref{eq:finite-propagation}.
\end{proof}

For fixed~$t$ and local~$f$, Lemma~\ref{lem:finite-propagation} implies uniform
convergence of the zero-boundary finite-volume semigroups to~$P_t f$ along any
exhaustion whose boundary recedes from the dependence set of~$f$. In
particular, the nonnegative centered-monomial coefficients of the finite-volume
semigroups used in Theorem~\ref{thm:centered-moment-order} pass to the
infinite-volume limit.
